\documentclass[11pt]{article}
\usepackage[a4paper,margin=1in]{geometry}
\usepackage{amsmath,amssymb,mathtools,bm}
\usepackage{enumitem,booktabs,array}
\usepackage{tcolorbox}
\usepackage{hyperref}
\usepackage[protrusion=true,expansion=false]{microtype}
\hypersetup{colorlinks=true,linkcolor=blue,urlcolor=blue,citecolor=blue}
\setlist[itemize]{topsep=2pt,itemsep=2pt}
\setlist[enumerate]{topsep=2pt,itemsep=2pt}
\newtcolorbox{keybox}{colback=blue!3!white,colframe=blue!40!black,boxrule=0.4pt,arc=1mm}
\newtcolorbox{alertbox}{colback=red!3!white,colframe=red!40!black,boxrule=0.4pt,arc=1mm}
\newtcolorbox{answerbox}{colback=green!3!white,colframe=green!40!black,boxrule=0.4pt,arc=1mm}
\newtcolorbox{drillbox}{colback=yellow!5!white,colframe=orange!60!black,boxrule=0.4pt,arc=1mm}
\newcommand{\EV}{\mathbb{E}}
\newcommand{\blank}[1]{\underline{\hspace{#1}}}

\title{W01-02: Gatev \& Strahan (2006, JF) \\
\large ``Banks' Advantage in Hedging Liquidity Risk'' --- Active Recall Workbook}
\author{Banking \& Risk Management --- Exam Prep}
\date{\today}
\begin{document}\maketitle

\begin{keybox}
\textbf{Paper in one sentence.} Empirical confirmation of KRS: when commercial-paper (CP) spreads widen, \emph{corporate firms draw their bank credit lines and simultaneously deposits flow into banks}, giving banks a natural hedge. The deposit-to-commitment correlation is \emph{negative} in market-stress states --- exactly the negative-$\rho$ condition KRS need.

\textbf{Placement.} Empirical companion to W01-01 (KRS). Motivates later crisis-era work (Ivashina--Scharfstein W04-04, Drechsler--Savov--Schnabl W02-03).
\end{keybox}

\section*{Round 1: Complete Walk-Through}

\subsection*{1.1 Research design}
\begin{itemize}
\item Sample: US bank holding companies, 1990--2002.
\item Key stress variable: \emph{CP spread} $=$ (commercial-paper rate $-$ T-bill rate). A widening CP spread signals that firms are losing access to CP markets.
\item Outcome variables at the bank level: (a) total deposits, (b) unused commitments drawn, (c) liquid-asset holdings.
\item Design is essentially a panel of bank-week or bank-quarter observations with time-series variation in the stress indicator.
\end{itemize}

\subsection*{1.2 Testable predictions from KRS}
\begin{enumerate}
\item When CP spreads widen (firms locked out of CP market), \emph{commitment drawdowns} rise.
\item At the same time \emph{deposit inflows} rise (flight to quality).
\item The two movements offset, so bank \emph{liquid assets} are relatively stable even in stress.
\item Banks with transaction-deposit-heavy bases receive the largest deposit inflows.
\end{enumerate}

\subsection*{1.3 Headline results}
\begin{itemize}
\item Strong, statistically significant negative correlation between deposit growth and drawdowns of credit lines in high-CP-spread weeks.
\item During the 1998 Russia/LTCM crisis: CP spreads spiked; firms tapped credit lines; deposits surged into banks, especially large banks with retail branch networks.
\item Transaction deposits explain the hedging: banks with high \emph{demand-deposit} shares show the strongest inflows under stress.
\item Liquid-asset buffers at banks remain stable --- consistent with internal offsetting.
\end{itemize}

\subsection*{1.4 Specification}
Panel regression of the form
\[
y_{bt} = \alpha_b + \delta_t + \beta\cdot \text{CP\_spread}_t + \gamma\cdot (\text{CP\_spread}_t\times \text{DepositMix}_b) + \varepsilon_{bt},
\]
with $y_{bt}$ alternately deposits, drawn commitments, or liquid assets. The interaction term identifies which banks capture the flight-to-quality deposits.

\subsection*{1.5 Identification \& econometric concerns}
\begin{alertbox}
\begin{itemize}
\item CP spread could reflect \emph{demand} shocks (firms doing worse) or \emph{supply} shocks (investors pulling out). GS argue either interpretation works for the deposit-hedge story.
\item Deposit growth at large banks could be driven by branching, M\&A, or other time trends --- hence bank fixed effects ($\alpha_b$) and time fixed effects ($\delta_t$).
\item External validity: 1990--2002 is pre-GFC. The Ivashina--Scharfstein (2010) and Acharya--Mora (2015) papers reassess this mechanism during 2008, with qualified support.
\end{itemize}
\end{alertbox}

\section*{Round 2: Level 1 Fill-in}
\begin{enumerate}
\item Stress indicator: \blank{2cm} spread $=$ CP rate $-$ \blank{1cm} rate.
\item In stress, firms draw \blank{2cm} and deposits flow \blank{1cm} banks.
\item Banks with more \blank{2cm} deposits see the largest inflows.
\item Bank \blank{2cm} assets remain stable even in high-stress weeks.
\end{enumerate}

\section*{Round 3: Level 2 Prompts}
\begin{enumerate}
\item Write GS's headline panel specification and explain every fixed effect.
\item Why does flight-to-quality generate negative $\rho$ between $\tilde W$ (deposit outflows) and $\tilde D$ (commitment drawdowns)?
\item How would you reproduce GS using 2008--09 crisis data? Which results strengthen, which weaken?
\item What role does deposit insurance play in this mechanism?
\end{enumerate}

\section*{Round 4: Minimal Scaffolding}
From a blank page: (i) sample and stress variable, (ii) four KRS predictions being tested, (iii) what is found, (iv) spec, (v) two external-validity concerns.

\section*{Round 5: Blank Page}
Why is GS the natural empirical follow-up to KRS? Write an essay linking the two.

\vspace{6pt}
\begin{drillbox}
\textbf{Speed Drill.} (1) Definition of CP spread. (2) Why deposits flow \emph{in} when firms draw lines. (3) Which bank type benefits most. (4) One sentence on the 1998 LTCM episode. (5) Why liquid assets stay stable.
\end{drillbox}

\section*{Quick Reference \& Answer Key}
\begin{answerbox}
\begin{itemize}
\item \textbf{CP spread.} Commercial paper rate minus T-bill rate; rises when investors leave the CP market.
\item \textbf{Flight-to-quality.} Investors pulling out of CP deposit with banks for safety; simultaneously firms that relied on CP draw their bank credit lines.
\item \textbf{Negative $\rho$.} Deposit inflow is positive \emph{exactly} when drawdowns are positive, giving internal offsetting $\Rightarrow$ $\rho<0$ in bad states.
\item \textbf{Cross-section.} Large, transaction-deposit banks are the ``liquidity providers of last resort'' for the non-bank corporate sector.
\end{itemize}
\end{answerbox}

\begin{keybox}
\textbf{30-second exam pitch.} Gatev--Strahan (2006) is the empirical proof of KRS. Using weekly bank-level data 1990--2002 and the commercial-paper spread as the stress variable, they show that when firms lose access to CP markets they draw their bank credit lines --- and at the same moment deposits flow \emph{into} banks, especially large banks with transaction-deposit bases. The two flows offset, so bank liquid assets stay stable. Economically, $\rho(\tilde W,\tilde D)<0$ in stress states. The paper explains why the modern US banking system is dominated by large, diversified, transaction-deposit-funded institutions: they are the natural liquidity insurer for the rest of the economy. Caveats: pre-GFC sample; later work (Ivashina--Scharfstein 2010, Acharya--Mora 2015) shows the hedge partially breaks down when the shock hits the banking system itself.
\end{keybox}

\end{document}
